% !TEX root = ../../cookbook.tex
\newpage
\recipe{The Perfect Cheesecake}{This recipe is reproduced almost verbatim from the version by Jerry P. Fairley's cookbook, in turn modified from the original found at \href{https://www.simplyrecipes.com/recipes/perfect_cheesecake/}{www.simplyrecipes.com}.}{}{}

\begin{multicols}{2}
	\subsection*{Ingredients}
	
	\subsubsection*{Crust}
	\begin{ingredients}
		\item 2 cups (200 g) graham cracker crumbs
		\item 2 tbsp (25 g) granulated sugar
		\item Pinch of salt
		\item 5 tbsp (70 g) unsalted butter
	\end{ingredients}
	
	\subsubsection*{Filling}
	\begin{ingredients}
		\item 2 lbs (900 g) cream cheese \textbf{at room T°}
		\item 1\sfrac{1}{3} cups (265 g) granulated sugar
		\item Pinch of salt
		\item 2 tsp vanilla extract
		\item 4 large eggs\footnote{US standard, $\sim$57 g each}
		\item \sfrac{2}{3} cup (160 g) sour cream
		\item \sfrac{2}{3} cup (160 g) heavy whipping cream
	\end{ingredients}
	\columnbreak
	\subsection*{ }
	\subsubsection*{Toppings \normalfont(optional)}
	\begin{ingredients}
		\item 2 cups (460 g) sour cream
		\item \sfrac{1}{3} cup (40 g) powdered sugar
		\item 1 tsp vanilla extract
		\item 12 oz (340 g) fresh raspberries, blackberries, or strawberries, finely diced
		\item \sfrac{1}{2} cup (100 g) granulated sugar
		\item \sfrac{1}{2} cup (120 g) water
	\end{ingredients}
	
	\subsubsection*{Special equipment}
	\begin{ingredients}
		\item a 9-in (23 cm) diameter, 2\sfrac{3}{4} inch (7 cm) high springform pan
		\item Heavy-duty \textbf{18-in (45 cm) wide} aluminum foil (\underline{do not} try to use regular-width aluminum foil)
		\item Parchment paper
		\item A large, high-sided roasting pan (large enough to contain the springform pan)
	\end{ingredients}
	
\end{multicols}

\medskip

% --- Preparation ---
\noindent \textbf{Crust:} \\ 
Unroll a sheet of \textbf{parchment paper} the size of the \textbf{springform pan}. 
Using the bottom of the springform pan as a pattern, draw a circle on the parchment paper and carefully cut it out. 
It should just fit in the bottom of the pan. 
Set the parchment circle aside for now.

Spread out a kitchen towel on the counter and set the springform pan, upside-down, on the towel\footnote{To avoid tearing or puncturing the delicate aluminum foil}.
Carefully and gently tear-off a square of the wide \textbf{aluminum foil} that will completely cover the entire bottom and sides of the springform pan. 
Place the foil over the bottom of the pan and gently fold the foil down over the sides, then turn the pan over and crimp the foil all around the lip. 
Do all of this very carefully, on top of the towel, so you don’t get even a little pinhole in the foil\textemdash if you do, start over with a new piece of aluminum foil. 
Once you get the first layer of aluminum foil on the springform pan, turn the pan upside-down again and put another piece of aluminum foil over the bottom and sides of the pan in exactly the same way as the first piece. 
This second layer of foil is an insurance policy in case the first layer has a tiny pinhole that you missed.

Put the \textbf{butter} in a microwave-proof bowl and heat it for 30 seconds. 
Wait for a couple of minutes, then swirl the bowl around and heat it again for 30 s. 
Let the bowl sit until any solid butter that is left melts completely.
Place a baking rack in the bottom third of the oven and preheat to 350°F (177°C). 
Roughly break-up the \textbf{graham crackers} and pulse them in a food processor until they are finely-ground. 
Put the crumbs in a small or medium mixing bowl. 
Stir in the \textbf{sugar} and pinch of \textbf{salt}, then pour in the melted butter and mix well with your hands.

Keeping the pan on the kitchen towel (to minimize the chance of pinholes in the foil), put the circle of parchment paper into the bottom, inside the foil-wrapped springform pan, and dump in all the graham cracker crust mixture. 
Press the crumbs into the bottom of the pan\footnote{I usually use the bottom of a measuring cup first, then my fingers (or a spoon) to finish the job.}, leaving just a slightly thicker crust along the outside edge of the pan. 
Be very careful not to tear the aluminum foil while you do this (or while doing anything else)\footnote{If it seems as though I’m making a really big deal out of getting a hole in the aluminium foil, that’s because I am. Even the tiniest hole will let water in the pan and spoil the crust. Making a good filling is tedious, but easy. The hardest part of making a good cheesecake is keeping the crust from being soggy. Trust me on this\textemdash it took me years to work out the crust part.}. 
Carefully (so you don’t make a hole in the aluminum foil) place the springform pan on the rack in the oven and bake for 10 minutes. 
Remove the pan and place it on the kitchen towel. Reduce the oven temperature to 325°F (325°C). \\ 

\noindent \textbf{Filling:} \\ 
You can do all the following steps with a handheld electric mixer; however, if you have one, it is (a lot) easier to use a stand-mixer, and the finished product is actually better. 
Although it’s hard work to mix the ingredients enough with a hand-mixer, with persisence, you can do it. 
The moral of the story is: if you want a good cheesecake, don’t short the mixing times below, just because you are using a hand-mixer! 

Put 2 quarts of water on to boil while you are preparing the filling, and keep them hot. 
You will need them after you make the filling, when you bake the cheesecake.

Cut the \textbf{cream cheese} up into chunks and place it all in a large mixing bowl. 
Mix the cream cheese\footnote{If you have a stand-mixer, use the paddle attachment.} on medium speed for at least 4 minutes, or until it is light and fluffy. 
Add the \textbf{sugar}, then beat the mixture for another 4 minutes. 
Add the \textbf{salt} and beat for 2 minutes, then add the \textbf{vanilla} and beat for another 2 minutes. 
Add the \textbf{eggs}, one at a time, beating the batter for 1 minute after adding each egg.
Add the \textbf{sour cream}, beating for 4 minutes, then add the \textbf{heavy cream}, beating for another 4 minutes. 
After each addition (sugar, salt, vanilla, eggs, sour cream, heavy cream), use a spatula to scrape the bottom and sides of the mixing bowl, to make sure all the little chunks of cream cheese get whipped and all the ingredients are completely incorporated. \\ 

\noindent \textbf{Baking the cheesecake:} \\ 
Carefully set the foil-wrapped springform pan in the large, high-sided roasting pan. 
Pour the cream cheese filling into the springform pan, on top of the graham cracker crust, and gently smooth the top with a rubber spatula if needed. 
Place the roasting pan (with the springform pan inside) on the lower rack (i.e., in the bottom third) of the oven. 
Pour the hot water (heated up in the previous step) into the roasting pan until it comes about halfway up the sides of the springform pan. 
The water should be 1\sfrac{1}{4} to 1\sfrac{1}{2} inches (3 to 4 cm) deep. 

Bake the cheesecake in the oven at 325°F (165°C) for at least 1\sfrac{1}{2} hours. 
After 1\sfrac{1}{2} hours, check to see if the cake is done by putting a bamboo skewer or a knife into the middle of the cake; if it comes out clean, the cake is done, but if there is goop (filling) on the knife, it needs more time. 
The top should be lightly browned; if it starts to look too brown, but the cheesecake still isn’t completely done, use a piece of aluminum foil to loosely cover the top. 
The exact baking time is probably highly dependent on the oven. 
The original recipe says 1\sfrac{1}{2} hours, but I have found it takes closer to 2\sfrac{1}{2} hours in our oven, and I sometimes resort to turning up the temperature to 350°F (175°C).
Make sure to top-up the water in the roasting pan if it gets too low, but I rarely find I need to add more water.

Once the cheesecake is set, turn off the oven and open the door about an inch, leaving the cheesecake inside. 
Let the cake cool for at least another hour, or longer if you like. 
I usually leave it until it is almost completely at room temperature. 
This helps keep the top of the cheesecake from developing cracks. 

Once the cheesecake is cooled, take it out of the oven and cover it with foil, without letting the foil touch the top of the cheesecake. 
Put it in the refrigerator and chill overnight, or for a minimum of 4 hours. \\

\noindent \textbf{Preparing the toppings:} \\ 
While the cheesecake is baking, place the \textbf{sour cream} in a medium sized bowl and stir in the powdered sugar and vanilla. 
Mix until smooth, then cover and chill until you are ready to serve the cake.

Put the \textbf{raspberries} (or \textbf{blackberries}, or \textbf{strawberries}, or whatever you are using), \textbf{sugar}, and \textbf{water} into a small saucepan. 
Simmer the mixture, occasionally using a potato masher to crush-up the fruit. 
After about 5–10 minutes the syrup will begin to get a little thicker. 
Once it starts to thicken, remove it from the heat and let it cool. 
Be careful not to overdo it, because the sauce will thicken more as it cools. \\

\noindent \textbf{Serving the Cheesecake:} \\ 
Take the cheesecake out of the refrigerator and remove the aluminum foil. 
Getting the springform pan off the cake without messing it up can be a little tricky. 
Try running a blunt knife around the inside of the pan, between the pan and the cake, before opening the latch on the pan. 
You can also try warming the sides of the pan (the original recipe suggests using a hair-dryer, but I’ve never had much success with that\textemdash I just use a knife). 
It might help to let the pan sit outside the refrigerator for 10–15 minutes before trying to unmold it, so the sides of the pan can warm up a little.

Anyway, once you unmold the springform pan from the cake, spread the top of the cheesecake with about half the sour cream topping. 
I suggest using a hot knife to cut the cake. 
I keep a kitchen towel (or paper towels) and a tall drinking glass full of hot water next to the cake; in between cutting each slice, I stick the knife in the water for a minute, then dry it off with the towel. 
This makes a mess of the towel, but you can always wash that later. 
Drizzle a little of the raspberry syrup on each slice and serve.

Note that the recipe for the sour cream topping makes more than you really need. 
I serve it on the side for people who want extra, but you can also serve the cheesecake without any sour cream dressing (or raspberry sauce, for that matter). 
The toppings are optional, and you can decide if you want to make (and use) them or not, depending on your tastes.

Making a cheesecake is a lot of work, although it does get easier after you’ve done it a few times. As I mentioned above, using a stand-mixer is a big help. Most of all, however,  
I have found the best cheesecakes are made with love! 
Love of baking, certainly, and also love for the person for whom you are making it. 
I don’t know if you’ve ever noticed this, but I never make a cheesecake for myself\textemdash it’s always easier to make it for someone that is dear to you.

\vspace*{0.75 cm}

% --- Description ---
\begin{recipeblock}
	\noindent
\end{recipeblock}
